% Vietnamese translation for OI Public Library
% Source-PDF: IOI中国国家候选队论文/2004/黄源河.pdf
\documentclass[11pt]{article}
\usepackage{oipl-vn}

\title{Bàn về xây dựng và ứng dụng mô hình đồ thị}
\author{Huang Yuanhe}
\date{}

\begin{document}
\maketitle

\origtitle{A brief discussion of building and applying graph-theory models}
\sourcepath{IOI China candidate team papers / 2004 / Huang Yuanhe.pdf}

\section*{Từ khóa}

Mô hình đồ thị; xây dựng; chuyển hóa.

\section*{Tóm tắt}

Trong những năm gần đây, các bài toán đồ thị xuất hiện ngày càng nhiều trong
các kỳ thi tin học. Là một nhánh toán học còn khá trẻ, lý thuyết đồ thị có
nhiều đặc điểm riêng so với các nhánh toán học khác. Khi dùng lý thuyết đồ
thị để giải bài, ta thường đạt được lời giải hiệu quả và gọn gàng. Tuy nhiên,
có công cụ này không có nghĩa là ta tự động giải được bài toán; điều quan
trọng còn nằm ở khả năng nhận ra và xây dựng thành thạo một loạt mô hình đồ
thị. Thông qua một số ví dụ, bài viết này giới thiệu ngắn gọn vài phương pháp
mô hình hóa bằng đồ thị.

\section{Mở đầu}

Khi dùng kiến thức toán học để giải bài, bước đầu tiên là phân tích bài toán
thực tế và xây dựng một mô hình toán học mô tả bài toán đó. Sau đó ta dùng lý
thuyết và phương pháp toán học để phân tích mô hình, thu được kết quả, rồi
quay lại giải quyết bài toán thực tế ban đầu. Mô hình đồ thị là một loại mô
hình toán học đặc biệt. Xây dựng mô hình đồ thị nghĩa là rút ra từ nguyên mẫu
của bài toán những thông tin và yếu tố hữu ích, rồi trừu tượng hóa bài toán
thành quan hệ giữa đỉnh, cạnh và trọng số.

Sau khi mô hình hóa bằng đồ thị, những thông tin rời rạc và hỗn độn trở nên
có quy luật; liên hệ nội tại giữa các yếu tố được thể hiện bằng quan hệ giữa
đỉnh, cạnh và trọng số. Nhiều mô hình đồ thị kinh điển có thể được giải trực
tiếp bằng những thuật toán xác định. Với một số bài toán phức tạp, chỉ cần
nhận rõ bản chất của bài toán, nắm được điểm then chốt và xây dựng một mô hình
đồ thị phù hợp, ta thường có thể chuyển chúng thành những bài toán kinh điển
quen thuộc.

Điều tôi muốn trình bày trong bài viết này chính là một vài kinh nghiệm và
nhận thức của bản thân về mô hình hóa bằng đồ thị.

\section{Phân tích ví dụ}

\subsection{Ví dụ 1: Place the Robots (ZOJ)}

\subsubsection*{Tóm tắt đề bài}

Cho một bàn cờ $N\times M$ với $N,M\le 50$. Mỗi ô thuộc một trong ba loại:
đất trống, cỏ, hoặc tường. Robot chỉ được đặt trên đất trống. Nếu hai robot
cùng hàng hoặc cùng cột và giữa chúng không có tường, chúng có thể tấn công
lẫn nhau. Hỏi với bàn cờ đã cho, có thể đặt nhiều nhất bao nhiêu robot sao
cho không robot nào tấn công robot nào?

\subsubsection*{Phân tích}

Trong nguyên mẫu của bài toán, các thông tin như cỏ và tường không phải là
điều ta quan tâm trực tiếp; ta chỉ quan tâm đến các ô trống và quan hệ giữa
chúng. Vì vậy, một mô hình đơn giản rất tự nhiên là: lấy mỗi ô trống làm một
đỉnh, nối cạnh giữa hai ô trống xung đột với nhau. Khi đó bài toán chuyển
thành tìm tập độc lập lớn nhất của đồ thị. Như đã biết, đây là một bài toán
NP-đầy đủ. Rõ ràng mô hình này không đem lại thuận lợi nào cho việc giải bài,
vì vậy ta phải xây dựng một mô hình mới thực sự hiệu quả.

Ta gọi một đoạn liên tiếp trong cùng hàng hoặc cùng cột, bị ngăn cách bởi
tường và có chứa ô trống, là một \term{khối}. Hiển nhiên trong mỗi khối nhiều
nhất chỉ đặt được một robot. Ta đánh số các khối theo hàng và theo cột riêng
biệt. Cụ thể, khối ngang sẽ thuộc phía $X$, khối dọc sẽ thuộc phía $Y$. Nếu
hai khối có chung một ô trống, ta nối một cạnh giữa chúng. Như vậy ta thu được
một đồ thị hai phía.

Trong mô hình này, mỗi cạnh biểu diễn một ô trống. Hai ô trống xung đột chắc
chắn tương ứng với hai cạnh có chung một đầu mút, vì chúng nằm trong cùng một
khối ngang hoặc cùng một khối dọc. Do đó, bài toán chuyển thành bài toán tìm
matching lớn nhất trên đồ thị hai phía. Đây là bài toán kinh điển trong lý
thuyết đồ thị và có thể giải bằng thuật toán Hungary.

\subsubsection*{Tiểu kết}

So sánh hai mô hình trên, mô hình thứ nhất quá đơn giản và không nhận ra bản
chất của bài toán; mô hình thứ hai nắm được đầy đủ liên hệ nội tại của bài
toán và xây dựng khéo léo một mô hình đồ thị hai phía. Vì sao lại có hai kết
quả khác hẳn nhau như vậy? Thứ nhất, góc nhìn phân tích bài toán khác nhau:
mô hình thứ nhất lấy ô trống làm đỉnh, còn mô hình thứ hai lấy ô trống làm
cạnh. Thứ hai, mô hình thứ nhất chọn chưa đủ thông tin từ nguyên mẫu, khiến
mô hình được xây dựng không giữ lại được tính chất quan trọng của nguyên mẫu;
trong khi đó mô hình thứ hai giữ được tính chất ``bàn cờ'' quan trọng. Từ đây
có thể thấy việc lựa chọn thông tin là một bước cực kỳ quan trọng trong mô
hình hóa bằng đồ thị.

\subsection{Ví dụ 2: Thuê thu ngân (ACM Tehran 2000)}

\subsubsection*{Mô tả bài toán}

Một siêu thị ở Tehran mở cửa 24 giờ mỗi ngày và cần một nhóm thu ngân để đáp
ứng nhu cầu phục vụ. Quản lý siêu thị thuê bạn giúp giải quyết bài toán này:
ở những thời đoạn khác nhau trong ngày, siêu thị cần số lượng thu ngân khác
nhau, chẳng hạn lúc nửa đêm chỉ cần một nhóm nhỏ còn buổi chiều thì cần rất
nhiều người. Mục tiêu của ông ta là thuê ít thu ngân nhất.

Quản lý đã cung cấp số thu ngân tối thiểu cần có trong từng giờ của một ngày:
$R_0,R_1,\ldots,R_{23}$. Trong đó $R_0$ là số thu ngân tối thiểu cần có từ
nửa đêm đến 1 giờ sáng, $R_1$ là số cần có từ 1 giờ đến 2 giờ sáng, v.v. Các
dữ liệu này giống nhau ở mọi ngày. Có $N$ người xin việc. Mỗi ứng viên $I$
trong 24 giờ sẽ bắt đầu làm từ một thời điểm cố định và làm liên tục đúng
8 giờ. Ký hiệu $t_I$ với $0\le t_I\le 23$ là thời điểm bắt đầu đó. Nghĩa là
nếu ứng viên thứ $I$ được nhận, người đó sẽ làm liên tục 8 giờ kể từ thời
điểm $t_I$.

Bạn cần viết chương trình, nhập $R_I$ với $I=0,\ldots,23$ và $t_I$ với
$I=1,\ldots,N$, tất cả đều là số nguyên không âm, rồi tính số thu ngân ít
nhất cần thuê để thỏa mãn các ràng buộc trên. Ở mỗi thời điểm, số thu ngân
đang làm việc được phép lớn hơn giá trị $R_I$ tương ứng.

\paragraph{Input.}
Dòng đầu của input là số bộ dữ liệu kiểm tra, không quá $20$. Trong mỗi bộ,
dòng đầu gồm 24 số nguyên biểu diễn $R_0,R_1,\ldots,R_{23}$ với $R_I\le
1000$. Dòng tiếp theo là $N$, số ứng viên, với $0\le N\le 1000$. Mỗi dòng
tiếp theo chứa một số nguyên $t_I$ với $0\le t_I\le 23$. Không có dòng trống
giữa hai bộ dữ liệu.

\paragraph{Output.}
Với mỗi bộ dữ liệu, xuất đúng một dòng chứa một số nguyên, là số thu ngân ít
nhất cần thuê. Nếu vô nghiệm, hãy xuất \texttt{No Solution}.

\subsubsection*{Phân tích}

Thoạt nhìn bài này rất dễ khiến ta nghĩ đến tham lam, quy hoạch động hoặc
mạng luồng, nhưng các thuật toán đó đều không xử lý được bài toán này một
cách trực tiếp.

Vì bài toán có nhiều ràng buộc, để làm rõ ý tưởng ta trước hết biểu diễn các
ràng buộc bằng ngôn ngữ toán học. Gọi $S[i]$ là tổng số thu ngân được thuê
trong các thời điểm từ $0$ đến $i$, và $W_i$ là số ứng viên bắt đầu làm việc
tại thời điểm $i$. Khi đó các ràng buộc trong đề bài có thể chuyển thành hệ
bất đẳng thức sau:
\[
\begin{cases}
0\le S[i]-S[i-1]\le W_i, & 0\le i\le 23,\\
S[i]-S[i-8]\le R_i, & 8\le i\le 23,\\
S[23]+S[i]-S[i+16]\le R_i, & 0\le i\le 7.
\end{cases}
\]

Hệ bất đẳng thức như vậy khiến ta nghĩ ngay đến hệ ràng buộc sai phân. Với
mỗi bất đẳng thức $S[i]-S[j]\le K$, ta thêm một cung có hướng từ đỉnh $j$ đến
đỉnh $i$ với trọng số $K$. Giá trị nhỏ nhất của $S[23]$ mà ta cần tìm sẽ là
độ dài đường đi ngắn nhất từ đỉnh $0$ đến đỉnh $23$. Tuy nhiên, cần chú ý bất
đẳng thức thứ ba ở trên: nó chứa ba ẩn, nên không thể biểu diễn trong đồ thị
như một quan hệ cạnh.

Đến đây tưởng như ta đã rơi vào bế tắc. Phải chăng bài này không thể giải
bằng hệ ràng buộc sai phân? Không, ta còn cần thêm một bước chuyển hóa. Nếu
xem $S[23]$ là ẩn thì chắc chắn không làm tiếp được. Nhưng nếu xem $S[23]$ là
một giá trị đã biết, bất đẳng thức thứ ba chỉ còn hai ẩn $S[i]$ và $S[i+16]$;
khi đó ta thêm cung từ đỉnh $i+16$ đến đỉnh $i$ với $0\le i\le 7$ và trọng số
$R_i-S[23]$. Như vậy, toàn bộ hệ bất đẳng thức có thể chuyển thành một đồ thị
có hướng. Nghiệm của ẩn $S[i]$ chính là đường đi ngắn nhất từ đỉnh $0$ đến
đỉnh $i$ trong đồ thị. Khi trong đồ thị tồn tại chu trình âm, hệ bất đẳng thức
vô nghiệm.

Cách làm trên xem $S[23]$ là giá trị đã biết, nhưng trên thực tế $S[23]$ không
những là ẩn mà còn chính là giá trị ta cần tối ưu. Ta xử lý bằng cách dùng
tìm kiếm nhị phân để thử giá trị của $S[23]$, dần thu hẹp phạm vi, và dùng
phương pháp lặp để kiểm tra có tồn tại chu trình âm hay không, tức kiểm tra
tính khả thi. Nếu ngay cả khi $S[23]=N$ vẫn không khả thi, ta xuất
\texttt{No Solution}; ngược lại ta xuất giá trị nhỏ nhất của $S[23]$. Độ phức
tạp thời gian là
\[
  O(24^3\log_2 N).
\]

\subsubsection*{Tiểu kết}

Bài này dùng lý thuyết hệ ràng buộc sai phân. Trong thi đấu, những hệ như vậy
không quá thường gặp, nhưng chúng có thể giải quyết khéo léo một số bài khó.
Mô hình của loại bài này thường không lộ rõ, cần có suy nghĩ và chuyển hóa.
Điểm then chốt khi làm dạng bài này là biểu diễn các ràng buộc trong đề bài
thành bất đẳng thức, rồi chuyển các bất đẳng thức đó thành mô hình đường đi
ngắn nhất hoặc dài nhất trên đồ thị.

\subsection{Ví dụ 3: Greedy Island (ZOJ)}

\subsubsection*{Tóm tắt đề bài}

Có $N$ thẻ, với $N\le 100000$. Mỗi thẻ có ba năng lực, giá trị lần lượt là
$A_i,B_i,C_i$. Mỗi thẻ có thể dùng một trong ba năng lực và mỗi thẻ chỉ được
dùng một lần. Ta cần chọn $A$ thẻ dùng năng lực thứ nhất, $B$ thẻ dùng năng
lực thứ hai, và $C$ thẻ dùng năng lực thứ ba, với $A+B+C\le 100$. Hãy tính
cần dùng những thẻ nào và dùng năng lực nào của từng thẻ để tổng giá trị năng
lực tương ứng là lớn nhất.

\subsubsection*{Phân tích}

Có nhiều cách giải bài toán tối ưu hóa, chẳng hạn quy hoạch động hoặc mạng
luồng. Bài này là một mô hình mạng luồng khá rõ.

Trong mô hình mạng luồng có nhiều loại trọng số: lưu lượng, sức chứa, và cũng
có thể có chi phí. Trong bài này, sức chứa có thể dùng làm ràng buộc lựa chọn
để bảo đảm tính hợp lệ của nghiệm; chi phí biểu diễn giá trị của lựa chọn để
bảo đảm tính tối ưu. Vì vậy, nói chính xác hơn, bài này là một mô hình luồng
cực đại chi phí cực đại.

Sau khi nhận rõ mô hình của bài toán, bước tiếp theo là dựng đồ thị:
\begin{itemize}
  \item Mỗi thẻ $i$ được biểu diễn bằng một đỉnh $i$. Thêm ba đỉnh
  $P_1,P_2,P_3$ biểu diễn ba loại năng lực, cùng với nguồn $S$ và đích $T$.
  \item Từ nguồn lần lượt nối cung đến $P_1,P_2,P_3$, có sức chứa tương ứng
  là $A,B,C$ và chi phí bằng $0$.
  \item Từ $P_1,P_2,P_3$ lần lượt nối cung đến mỗi đỉnh $i$ với $1\le i\le N$,
  sức chứa bằng $1$, chi phí lần lượt là $A_i,B_i,C_i$.
  \item Từ mỗi đỉnh $i$ với $1\le i\le N$ nối một cung đến đích, sức chứa
  bằng $1$ và chi phí bằng $0$.
\end{itemize}

Nói cách khác, mạng có dạng
\[
  S \to \{P_1,P_2,P_3\} \to \{1,2,\ldots,N\} \to T,
\]
trong đó các cung từ $S$ đến ba đỉnh năng lực quyết định số thẻ cần chọn cho
từng loại, còn các cung từ đỉnh năng lực đến thẻ mang giá trị thu được khi
dùng thẻ đó cho năng lực tương ứng.

Sau khi dựng đồ thị, ta tìm luồng cực đại chi phí cực đại từ $S$ đến $T$, rồi
kiểm tra các cung chảy ra từ $P_1,P_2,P_3$ để xuất phương án tối ưu. Thuật
toán này còn rất thô: độ phức tạp thời gian là $O(N^3)$, độ phức tạp bộ nhớ
là $O(N^2)$, nên cả thời gian lẫn bộ nhớ đều không khả thi và cần tiếp tục
tối ưu.

Bài này có đến một trăm nghìn thẻ, trong khi số thẻ cuối cùng cần chọn không
vượt quá $100$. Nếu trước khi dựng đồ thị ta xóa bớt các thẻ vô dụng, hiệu
quả chắc chắn sẽ tăng đáng kể.

Những thẻ nào là vô dụng?

Trước hết xét việc chọn năng lực thứ nhất. Nếu sắp xếp toàn bộ thẻ theo giá
trị năng lực thứ nhất giảm dần, hiển nhiên ta nên cố gắng chọn $A$ thẻ ở phía
đầu. Vì mỗi thẻ chỉ được dùng một lần, chúng có thể xung đột với hai loại
năng lực còn lại; số thẻ xung đột nhiều nhất là $B+C$. Do đó các thẻ thực sự
có ích đối với năng lực thứ nhất chỉ là $A+B+C$ thẻ đầu tiên.

Tương tự, với năng lực thứ hai và thứ ba, ta cũng chỉ cần giữ lại $A+B+C$ thẻ
có giá trị lớn nhất của từng năng lực. Bước này có thể thực hiện trong thời
gian tuyến tính.

Đây là một phương pháp vừa đơn giản vừa hiệu quả. Sau bước xử lý này, số thẻ
được giữ lại không vượt quá $3(A+B+C)$, số đỉnh giảm mạnh, và thời gian tìm
luồng cực đại chi phí cực đại giảm xuống còn
\[
  O((A+B+C)^3).
\]
Đến đây thuật toán đã được tối ưu đến mức chấp nhận được, với độ phức tạp
tổng thể chỉ là
\[
  O(N+(A+B+C)^3).
\]

Nếu muốn tăng hiệu quả hơn nữa, có thể dùng thuật toán tốt hơn để xóa các
đỉnh dư thừa. Tuy nhiên điều đó không có nhiều ý nghĩa và cũng không phải
trọng tâm thảo luận của bài viết này.

Ngoài ra, bài này còn có thể chuyển thành mô hình đồ thị hai phía và giải
bằng thuật toán matching tối ưu. Bước này xin để độc giả tự suy nghĩ.

\subsubsection*{Tiểu kết}

Bài này xây dựng mô hình mạng luồng. Hệ số trong các thuật toán thuộc lớp mô
hình này thường lớn, độ phức tạp lập trình cũng cao, nên trong thi đấu nó
thường được xem như một ``thuật toán dự phòng'' khi không còn hướng đi nào
khác. Tuy vậy, mô hình mạng luồng có phạm vi áp dụng rộng. Với một số bài
toán khá phức tạp hoặc có nhiều ràng buộc, mô hình mạng luồng có thể xử lý
rất tốt; hơn nữa những bài dựa trên mô hình mạng luồng thường tương đối dễ
nhận ra. Vì thế mô hình mạng luồng có ứng dụng rất rộng rãi.

\section{Tổng kết}

Bài toán biến hóa muôn hình vạn trạng, nên không có một quy tắc chung nào cho
việc xây dựng mô hình đồ thị của bài toán. Các ví dụ phía trên đều tương đối
đơn giản. Trong những bài toán phức tạp hơn, đôi khi ta cảm thấy khó xây dựng
một mô hình thích hợp. Khi đó, với tiền đề không làm thay đổi tính chất của
nguyên mẫu bài toán, ta cần trừu tượng hóa và đơn giản hóa nguyên mẫu, rồi
trên cơ sở đó xây dựng một mô hình phù hợp và hiệu quả. Đôi khi sau khi đã
xây dựng được một mô hình cho bài toán, ta lại cảm thấy mô hình đó khó giải.
Khi ấy ta có thể cần sửa đổi hoặc chuyển hóa mô hình, hoặc phân tích sâu hơn
nguyên mẫu bài toán, rút ra những yếu tố then chốt hơn để xây dựng một mô hình
dễ giải. Tất cả những việc này đều đòi hỏi kinh nghiệm phong phú, tư duy linh
hoạt và khả năng sáng tạo tốt.

\section*{Tài liệu tham khảo}

\begin{enumerate}
  \item Wu Wenhu, Wang Jiande, \emph{Phân tích thuật toán thực dụng và thiết
  kế chương trình}, Nhà xuất bản Công nghiệp Điện tử, 1998.
  \item Wu Wenhu, Wang Jiande, \emph{Hướng dẫn thi Olympic Tin học quốc tế và
  toàn quốc cho thanh thiếu niên: thuật toán và thiết kế chương trình trong lý
  thuyết đồ thị}, Nhà xuất bản Đại học Thanh Hoa, 1997.
  \item Bản thảo thảo luận giai đoạn 3 của đội tuyển tập huấn quốc gia IOI
  2003.
\end{enumerate}

\appendix

\section{Nguyên đề ví dụ 1}

Nguồn đề: ZOJ Monthly, October 2003 Contest.

\subsection*{Place the Robots}

Robert là một kỹ sư nổi tiếng. Một ngày nọ, anh được sếp giao một nhiệm vụ.
Bối cảnh của nhiệm vụ như sau.

Cho một bản đồ gồm các ô vuông. Có ba loại ô: tường, cỏ và ô trống. Sếp của
Robert muốn đặt nhiều robot nhất có thể trên bản đồ. Mỗi robot cầm một vũ khí
laser có thể bắn đồng thời theo bốn hướng bắc, đông, nam và tây. Robot phải
đứng yên tại ô được đặt ban đầu và luôn bắn. Tia laser chắc chắn có thể đi qua
ô cỏ, nhưng không thể đi qua ô tường. Robot chỉ được đặt trên ô trống. Dĩ
nhiên sếp không muốn thấy robot này làm hại robot khác. Nói cách khác, hai
robot không được nằm trên cùng một hàng hoặc cùng một cột, trừ khi giữa chúng
có tường.

Vì bạn là một lập trình viên thông minh và là một trong những người bạn tốt
nhất của Robert, anh ấy nhờ bạn giúp giải bài toán này: với mô tả của bản đồ,
hãy tính số robot lớn nhất có thể đặt trên bản đồ.

\paragraph{Input.}
Dòng đầu chứa một số nguyên $T$ với $T\le 11$, là số bộ dữ liệu.

Với mỗi bộ dữ liệu, dòng đầu chứa hai số nguyên $m,n$ với $1\le m,n\le 50$,
là số hàng và số cột của bản đồ. Sau đó là $m$ dòng, mỗi dòng chứa $n$ ký tự
\texttt{\#}, \texttt{*}, hoặc \texttt{o}, lần lượt biểu diễn tường, cỏ và ô
trống.

\paragraph{Output.}
Với mỗi bộ dữ liệu, trước hết xuất số thứ tự bộ dữ liệu trên một dòng theo
dạng \texttt{Case :id}, trong đó \texttt{id} là số thứ tự tính từ $1$. Dòng
thứ hai chỉ xuất số robot lớn nhất có thể đặt trên bản đồ đó.

\paragraph{Sample Input.}
\begin{verbatim}
2
4 4
o***
*###
oo#o
***o
4 4
#ooo
o#oo
oo#o
***#
\end{verbatim}

\paragraph{Sample Output.}
\begin{verbatim}
Case :1
3
Case :2
5
\end{verbatim}

\section{Nguyên đề ví dụ 3}

Đề có chỉnh sửa. Nguồn đề: ZOJ Sunny Cup 2003 Online Contest.

\subsection*{Greedy Island}

Gon và Killua đang tham gia một trò chơi tên là Greedy Island. Mục tiêu của
trò chơi là thu thập 100 thẻ phép thuật rải rác khắp trò chơi. Mỗi thẻ phép
thuật có ba thuộc tính: tấn công, phòng thủ và đặc biệt. Giá trị số của mỗi
thuộc tính nằm trong khoảng từ $0$ đến $100$. Mỗi thẻ chỉ được dùng một lần.
Tất cả thẻ phép thuật phải được lưu trong Book, vật phẩm ban đầu của trò chơi.
Book có thể chứa nhiều nhất 50 thẻ phép thuật, nên Gon và Killua tổng cộng có
thể có nhiều nhất 100 thẻ phép thuật. Hiện tại Gon và Killua có $n$ thẻ phép
thuật, và họ muốn dùng $A$ thẻ cho tấn công, $B$ thẻ cho phòng thủ, và $C$ thẻ
cho mục đích đặc biệt, với $A+B+C\le 100$. Nếu $n>A+B+C$, họ phải bỏ đi một
số thẻ.

Họ muốn biết tổng lớn nhất có thể của giá trị tấn công trong nhóm tấn công,
giá trị phòng thủ trong nhóm phòng thủ và giá trị đặc biệt trong nhóm đặc
biệt. Nếu có nhiều phương án, chọn phương án có tổng tất cả thuộc tính của
tất cả các thẻ được chọn là lớn nhất.

\paragraph{Input.}
Dòng đầu chứa một số nguyên $T$ với $1\le T\le 10$, là số bộ dữ liệu.

Với mỗi bộ dữ liệu, dòng đầu chứa một số nguyên $n$ với $n\le 100100$; dòng
tiếp theo chứa ba số nguyên $A,B,C$ với $A,B,C\ge 0$ và $A+B+C\le n$; tiếp đó
là $n$ dòng, mỗi dòng chứa giá trị tấn công, phòng thủ và đặc biệt của một thẻ
phép thuật.

\paragraph{Output.}
Với mỗi bộ dữ liệu, in trên một dòng hai số, cách nhau bởi một dấu cách: tổng
lớn nhất của giá trị tấn công trong nhóm tấn công, giá trị phòng thủ trong
nhóm phòng thủ và giá trị đặc biệt trong nhóm đặc biệt; và tổng lớn nhất của
tất cả thuộc tính của tất cả các thẻ.

\paragraph{Sample Input.}
\begin{verbatim}
2
3
1 1 1
100 0 0
0 100 0
0 0 100
4
1 1 1
12 32 44
33 48 37
37 38 33
46 79 78
\end{verbatim}

\paragraph{Sample Output.}
\begin{verbatim}
300 300
163 429
\end{verbatim}

\end{document}
